\section{Cost and Operating Economics}
\label{sec:economics}

\subsection{What a test costs}

Inference is the only input that scales with use, and it separates into two
components that behave quite differently. Both are measured against a
provider's billing record rather than derived from a price card.

\begin{table}[H]
  \centering\small
  \begin{tabularx}{\textwidth}{@{}>{\hsize=0.62\hsize}Y>{\hsize=1.66\hsize}Y>{\hsize=0.42\hsize}Y@{}}
    \toprule
    Item & Measured basis & Cost \\
    \midrule
    One executed test & 10.1 device steps across 232 stored observations; accessibility tree averaging 4{,}686 bytes per step; planner logging 384{,}000 tokens over 108 calls & USD 0.008 \\
    \addlinespace
    One 25-test campaign & Linear in tests & USD 0.20 \\
    \addlinespace
    One 500-test nightly suite & Linear in tests & USD 4 \\
    \addlinespace
    One specification ingestion & Roughly 22 calls on the most expensive tier, three self-consistency samples, no output cap & USD 1.80 \\
    \bottomrule
  \end{tabularx}
\end{table}

Within the per-test figure the player accounts for about 92\,\% and the planner
about 8\,\%. This is why reducing median steps from 50 to 9 was a real saving
and why optimising the planner prompt would not have been.

Ingestion behaves in the opposite way. It costs the equivalent of roughly two
hundred test executions, and it is paid once per document rather than once per
run, so it is a fixed cost rather than an operating one. It is also the largest
single line item in the project's own spending, because the specification was
re-read from scratch every time the ingestion pipeline changed during
development. Three metering controls would close that gap and are not yet
built: an output ceiling on extraction, a reduced sample count once extraction
has stabilised, and a cache that refuses to re-read an unchanged document.

\subsection{Comparison with published figures}

Published costs for comparable systems are USD 13 to 22 per application for a
two-hour budget, and USD 10.17 to 18.76 per thousand queries. At USD 0.008 per
executed test the system reported here is two to three orders of magnitude
below those figures. The comparison is indicative rather than like-for-like,
because the units differ and the tasks differ, but the gap is large enough to
survive a generous restatement of either side.

\subsection{What it would cost to self-host}

The project had no accelerator hardware, so hosted inference was the only route
available. The question of what the system would cost without that dependency
is worth answering with the arithmetic shown, because the answer is not the
obvious one.

The player dominates: it is invoked on every step, accounts for about 92\,\% of
the cost of a test, and the checkpoint used for the reported campaigns has 32
billion parameters. Weights at 16-bit precision need roughly 64\,GB, at 8-bit
roughly 32\,GB, and at 4-bit roughly 16 to 18\,GB, before the key-value cache,
the vision encoder's activations and one screenshot per step. The practical
minimum is a single 48\,GB accelerator at 8-bit, or two 24\,GB consumer cards at
4-bit.

\begin{table}[H]
  \centering\small
  \begin{tabularx}{\textwidth}{@{}>{\hsize=0.62\hsize}Y>{\hsize=1.60\hsize}Y>{\hsize=0.48\hsize}Y@{}}
    \toprule
    Route & Basis & Cost per test \\
    \midrule
    Hosted API, what we did & Measured against billing at 10.1 steps per test & USD 0.008 \\
    \addlinespace
    Rented accelerator & One 48\,GB cloud instance at about USD 1.25 per hour run continuously for a year, roughly USD 11{,}000, producing about 260{,}000 tests at one test per two minutes of occupied card & about USD 0.042 \\
    \addlinespace
    Purchased accelerator & A 48\,GB card at USD 5{,}000 to 8{,}000 in a host costing USD 8{,}000 to 12{,}000, amortised over three years, ignoring power and administration & about USD 0.013 \\
    \bottomrule
  \end{tabularx}
\end{table}

Prices are indicative at the time of writing and are the input a reader should
replace with their own. The structure of the comparison is the part that does
not change.

\subsection{Why renting a card costs more than renting tokens}

The result is counter-intuitive only until the duty cycle is examined. A test
occupies its device for up to 900 seconds but consumes accelerator time only in
the brief interval between observing a screen and deciding the next action, on
the order of a hundred seconds across 10.1 steps. The rest is the application
responding, animations completing and the emulator settling. A dedicated card
idles for most of a run and the tenant pays for that idling; a hosted provider
does not charge for it, because it fills those gaps with other tenants' work.

\keypoint{The card becomes economical exactly when it can be kept busy, which
means many devices driven in parallel against one served model. That is the
natural shape of a nightly regression fleet inside an organisation, and it was
not the shape of this project. At high utilisation a purchased card approaches
and then undercuts the API price.}

There is also an argument for self-hosting that does not depend on cost at all.
Every screenshot of the application under test is currently transmitted to a
third party. An organisation with a data-residency obligation, or one testing
an unreleased product, would need to self-host regardless of the arithmetic
above. Because the models in use are open-weight checkpoints rather than
proprietary endpoints, that path stays open, and keeping it open is why the
model provider has sat behind a swappable interface since June.
